% ARCHIVED FROM v82 EXACT SOURCE EXCERPT: 02_fixed-h-mcfg/main.tex
% Source reference: archive/mcfg-v82-pretrim-full-20261018 (commit bcc0caf0eade66f552127171c0de24d2b92950b0)
% Restoration: replace the corresponding trimmed location with this excerpt (without these header lines).

\begin{proposition}[Fixed windows do not rescue the bracketed CFG family at fan-out two]
\label{prop:fixed-window-doubling-fanout-two}
Let $\Gamma_{\mathrm{br}}=\{\mathtt a,\lfloor,\rfloor,\lceil,\rceil\}$,
let $L_n$ be generated by
\[
 E_0\to\mathtt a,\qquad
 E_i\to\lfloor E_{i-1}E_{i-1}\rfloor
       \mid\lceil E_{i-1}\rceil\quad(1\le i\le n),
\]
with start symbol $E_n$, as in the companion CFG
manuscript~\cite{kuriyamaCFG}.
Let $h_{\mathrm U}$ send $\lceil,\rceil$ to the absorbing
element of the two-element semilattice and all other letters
to its identity.
For any $k,\ell\ge0$ and
$n\ge\max\{3,k+\ell+1\}$,
\[
 L_n\notin\Cmcf{2}{h_{\Gamma_{\mathrm{br}}}^{k,\ell}},
 \qquad
 L_n\notin\Cmcf{2}{(h_{\Gamma_{\mathrm{br}}}^{k,\ell},h_{\mathrm U})}.
\]
Thus no \emph{fixed} prefix--suffix window, even combined with
$h_{\mathrm U}$, repairs this family for the higher-arity
substitutability required at fan-out two.
\end{proposition}
\begin{proof}
Define $W_0=\mathtt a$ and
$W_i=\lfloor W_{i-1}W_{i-1}\rfloor$.
The unique word in $L_n$ with no ceiling symbols is $W_n$.
Since $W_{n-1}$ is the complete binary-tree word, write
$W_{n-1}=B\mathtt a R$ with
$B=\lfloor^{n-1}$.
Then
\[
 W_n=\lfloor B\mathtt a R\,B\mathtt a R\rfloor.
\]
Both $B$ and $R$ are nonempty, with
$|B|=n-1\ge k+\ell$ and
$|R|=3\cdot2^{n-1}-n-2\ge n-1\ge k+\ell$.

Consider two tuples and two sentence contexts:
\[
 \begin{aligned}
 \tuple x&=(B,\ R B\mathtt a R),&
 \tuple y&=(B\mathtt a R B,\ R),\\
 E&=\lfloor\blank_1\mathtt a\blank_2\rfloor,&
 F&=\blank_1\lfloor\mathtt a\blank_2\rfloor.
 \end{aligned}
\]
Both contexts have a nonempty internal separator, as
required by Definition~\ref{def:tuple-substitutable}.
The first components $B$ and $B\mathtt a R B$ share
their length-$k$ prefix and length-$\ell$ suffix,
and the second components $R B\mathtt a R$ and $R$
do likewise. All four components have length at least
$k+\ell$ and contain no ceiling symbol.
Thus $\tuple x$ and $\tuple y$ have identical
componentwise values for both typings in the statement.
But
\[
 E[\tuple x]=E[\tuple y]=F[\tuple x]=W_n,
 \qquad
 F[\tuple y]=B\mathtt a R B\lfloor\mathtt a R\rfloor
 \ne W_n.
\]
Indeed, $F[\tuple y]$ has its first $\mathtt a$
after $n-1$ leading $\lfloor$ symbols, whereas
$W_n$ has $n$ such symbols.
As $F[\tuple y]$ also has no ceiling symbols,
it does not belong to $L_n$.
This contradicts $(2,h)$-tuple substitutability
for each specified typing.
\end{proof}

\begin{remark}
The proposition proves failure of \emph{class membership}
for the original CFG family, not a superpolynomial
characteristic-data lower bound at fan-out two.
In particular, it does not resolve
Open Problem~\ref{op:fanout-two}, which requires a
\emph{different} family actually belonging to
$\Cmcf{2}{h}$.
\end{remark}



\subsection{Rare Types and Two-Component Mobility}
\label{subsec:rare-type-mobility}

\begin{lemma}[Deleting a complete $h$-type]
\label{lem:mcfg-delete-type}
For any $f\ge1$, $L\in\Cmcf{f}{h}$, and $m\in M$,
the language $L^{\setminus m}:=L\setminus h^{-1}(m)$
also belongs to $\Cmcf{f}{h}$.
\end{lemma}
\begin{proof}
The fiber complement $h^{-1}(M\setminus\{m\})$ is
regular and $(f,h)$-tuple-substitutable: equal componentwise
types imply equal total types under every sentence context.
Intersections preserve $(f,h)$-tuple substitutability
(Proposition~\ref{prop:abelian-coset-filter}),
and $f$-MCFLs are closed under regular intersection
\cite{SekiEtAl1991}.
\end{proof}

\begin{proposition}[Rare types force long positive data]
\label{prop:mcfg-rare-type-sample}
Let $L\in\Cmcf{f}{h}$ and $m\in h(L)$, and put
$\mu_m(L)=\min\{|w|:w\in L,\ h(w)=m\}$.
Every characteristic sample for either
$\mathcal B^{\le r}_{f,h}$ ($r\ge1$) or
$\mathcal B_{f,h}$ contains a word of type $m$, and hence has
encoded total length at least $\mu_m(L)+1$.
More generally, the same necessity holds for any set-driven
reconstruction operator admitting characteristic samples for
both $L$ and $L^{\setminus m}$.
\end{proposition}
\begin{proof}
For a sample $K\subseteq L$ avoiding type $m$ we have
$K\subseteq L^{\setminus m}$.
Lemma~\ref{lem:mcfg-delete-type} and
Proposition~\ref{prop:reconstruction-soundness}, applied
with target $L^{\setminus m}$, show that both reconstructions
from $K$ are contained in this proper sublanguage of $L$.
For the general set-driven claim, suppose $\mathcal F$ is
characteristic for $L$ but avoids type $m$, and let
$\mathcal F'$ be characteristic for $L^{\setminus m}$.
The same finite set $\mathcal F\cup\mathcal F'$ belongs to
both target languages and contains both characteristic
samples, contradicting their distinct required outputs.
\end{proof}

\begin{lemma}[Power transfer at fan-out two]
\label{lem:mcfg-power-transfer}
Suppose $f\ge2$, $L\in\Cmcf{f}{h}$ and $u\in\Sigma^+$.
Let $g=h(u)$ and choose $i_0,\pi\ge1$ such that
$g^{i+\pi}=g^i$ for every $i\ge i_0$.
If $s,r\ge i_0$, $d\in\pi\mathbb Z_{>0}$, and some
word in $L$ contains $u^{s+d+r+2}$, then for
all $X,Z\in\Sigma^*$ and $Y\in\Sigma^+$,
\[
 Xu^sYu^{d+r}Z\in L
 \quad\Longleftrightarrow\quad
 Xu^{s+d}Yu^rZ\in L.
\]
\end{lemma}
\begin{proof}
Write the assumed positive word as $Pu^{s+d+r+2}Q$.
The tuples $\tuple x=(u^s,u^{d+r})$ and
$\tuple y=(u^{s+d},u^r)$ have equal componentwise
$h$-types.  Their common accepting sentence context
is $E=Pu\,\blank_1u\blank_2Q$: indeed
$E[\tuple x]=E[\tuple y]=Pu^{s+d+r+2}Q$.
Tuple substitutability therefore equates their
distributions.  Test these distributions in the
arbitrary valid context $X\blank_1Y\blank_2Z$.
\end{proof}

\begin{lemma}[Bounded local shift pair]
\label{lem:mcfg-local-shift}
Assume $\Sigma$ is nonempty. Let $c=|M|^2|\Sigma|$
and $R_c(5)$ be the
$c$-colour Ramsey number for $K_5$.
Every $w\in\Sigma^*$ with $|w|\ge R_c(5)$
admits a factorization
$w=P\,x\,p\,q\,p\,y\,Q$,
where $x,y\in\Sigma^+$, $p\in\Sigma$, $q\in\Sigma^*$,
\[
 h(x)=h(xpq),\qquad h(qpy)=h(y).
\]
Consequently, if $w\in L\in\Cmcf{2}{h}$, the tuples
$(x,qpy)$ and $(xpq,y)$ have the same distribution.
\end{lemma}
\begin{proof}
Colour every pair of cuts $i<j$ among the first
$R_c(5)$ cut positions by
\[
 (h(w[i:j]),\,h(w[i+1:j]),\,w[i]).
\]
A monochromatic five-set $i_1<\dots<i_5$ exists.
Choose $x=w[i_1:i_2]$, $p=w[i_2]=w[i_3]$,
$q=w[i_2+1:i_3]$,
$y=w[i_3+1:i_5]$, and take $P,Q$ to be
the prefix before $i_1$ and suffix after $i_5$.
The first colours of $(i_1,i_2)$ and $(i_1,i_3)$
give $h(x)=h(xpq)$; the second colours of
$(i_2,i_5)$ and $(i_3,i_5)$ give $h(qpy)=h(y)$.
The nonempty-separator context
$P\blank_1p\blank_2Q$ accepts both tuples,
so their distributions agree.
\end{proof}

\begin{theorem}[No fixed finite typing rescues the bracketed doubling family]
\label{thm:mcfg-all-typing-doubling-obstruction}
Let $h:\Gamma_{\mathrm{br}}^*\to M$ be an
\emph{arbitrary} finite-monoid morphism.
For the family $L_n$ of
Proposition~\ref{prop:fixed-window-doubling-fanout-two},
\[
 n\ge2|M|+3\quad\Longrightarrow\quad
 L_n\notin\Cmcf{2}{h}.
\]
The same conclusion holds for $\Cmcf{f}{h}$ whenever
$f\ge2$ and the representational fan-out is at least one.
\end{theorem}
\begin{proof}
Let $g=h(\lfloor)$, and choose an eventual index
$i_0\ge1$ and period $\pi\ge1$ for its powers,
with $i_0+\pi\le |M|+1$ and $i_0\le|M|$.
It suffices to consider $n\ge2i_0+\pi+2$.
Put $s=i_0$, $d=\pi$ and $r=n-s-d-2\ge i_0$.
The unique ceiling-free word
$W_n=\lfloor W_{n-1}W_{n-1}\rfloor\in L_n$
starts with $\lfloor^n$, so the hypothesis of
Lemma~\ref{lem:mcfg-power-transfer} holds with $u=\lfloor$.

Within this same $W_n$, take the first block
$\lfloor^s$ from its initial run, starting after
the very first $\lfloor$, and take the second
block $\lfloor^{d+r}$ from the initial run of
the \emph{second} copy of $W_{n-1}$.
The substring $Y$ between these blocks contains
the first $\mathtt a$ and is nonempty, so
$W_n=X\lfloor^sY\lfloor^{d+r}Z$ for suitable $X,Z$.
Power transfer would imply
$X\lfloor^{s+d}Y\lfloor^rZ\in L_n$.
This word has no ceiling brackets but has $n+d$
initial $\lfloor$ symbols before its first $\mathtt a$,
whereas the sole ceiling-free member $W_n$ has
exactly $n$.  This is a contradiction.
The uniform threshold follows from
$2i_0+\pi+2\le2|M|+3$.
\end{proof}

\begin{remark}
Theorem~\ref{thm:mcfg-all-typing-doubling-obstruction}
strengthens the restricted-window obstruction in
Proposition~\ref{prop:fixed-window-doubling-fanout-two},
whose sharper window-specific threshold is retained.
It is a \emph{nonmembership} theorem, not the
higher-fan-out sample-length lower bound sought in
Open Problem~\ref{op:fanout-two}.
\end{remark}

The fixed-window--finite-group extension is proved in
Appendix~\ref{app:window-group}; it applies to arbitrary
finite groups, with the sharper abelian constant retained
in Appendix~\ref{app:window-abelian}.
Whether the exponents two and three in the respective
window-only and window--finite-group characteristic-data bounds
are optimal is open, as is whether the doubly exponential
dependence on $f$ of the nonabelian constant $M_{A,f}$ can be reduced.
Appendix~\ref{app:local-group} extends this original-thickness
control to all typings whose positive image is a local group, the
largest class excluding the two-element semilattice.  The matching
lower bound is known only at fan-out one
(Corollary~\ref{cor:pseudovariety-dichotomy}).
Corollary~\ref{cor:mcfg-li-group-polydata} gives a
transparent sufficient condition on the positive-image semigroup in
$\mathbf{LI}\vee\mathbf{Grp}$.  The broader local-group extension and
pseudovariety-level dichotomy are treated separately in
Appendix~\ref{app:local-group}; higher-fan-out lower bounds remain open.

Other directions include adaptive learning over several fan-out biases,
learning useful typings without an a priori codomain bound, characterizing the
rank-one subfamily inside the typed classes, understanding the relation with
full-multicontext congruential MCFGs, and extending the reconstruction to
PMCFG-style copying.

\section*{Acknowledgments}
The author thanks Makoto Kanazawa and Ryo Yoshinaka for valuable discussions
and encouragement.

\section*{Declaration of competing interest}
The author declares that he has no known competing financial interests or
personal relationships that could have appeared to influence the work reported
in this paper.

\section*{Funding}
This research received no specific grant from any funding agency in the public,
commercial, or not-for-profit sectors.

\section*{Data availability}
No empirical datasets were used. The source manuscript and scripts
for reproducing finite illustrative checks are available in the
public repository
\url{https://github.com/growupkuriyama-hub/Papers/tree/main/02_fixed-h-mcfg}.
The general theorems are supported by mathematical proofs in this article,
not by exhaustive computational testing.

\section*{Declaration of generative AI and AI-assisted technologies in the manuscript preparation process}
The author used ChatGPT (OpenAI) and Claude (Anthropic) during the research
and manuscript-preparation process, including exploratory mathematical
argument development, proof drafting and criticism, literature and
consistency checks, and editing of text and LaTeX source.
AI-generated suggestions and reviews are not treated as independent
mathematical verification. The author takes responsibility for the
mathematical claims, citations, and final content of this article.

\appendix
